% Generated draft: requires mathematical and source review.
\documentclass[11pt]{article}
\usepackage[margin=25mm]{geometry}
\usepackage{fontspec}
\setmainfont{Latin Modern Roman}
\usepackage{amsmath,amssymb,mathtools,unicode-math}
\setmathfont{Latin Modern Math}
\usepackage{enumitem,needspace,xurl,hyperref}
\setlength{\parindent}{0pt}
\setlength{\parskip}{5pt}
\setlength{\emergencystretch}{3em}
\newcommand{\N}{\mathbb{N}}
\newcommand{\Z}{\mathbb{Z}}
\newcommand{\Q}{\mathbb{Q}}
\newcommand{\R}{\mathbb{R}}
\newcommand{\C}{\mathbb{C}}
\newcommand{\F}{\mathbb{F}}
\newcommand{\A}{\mathbb{A}}
\newcommand{\PP}{\mathbb{P}}
\newcommand{\E}{\mathbb{E}}
\newcommand{\eps}{\varepsilon}
\newcommand{\dd}{\,\mathrm{d}}
\begin{document}
\section*{UnsolvedMath ID: 20000453}
\textbf{Source problem code:} catalogue entry
\subsection{Definitions and mathematical statement}
Put $g(t)=3(t^4-54t^2-117t-243)$, and let $S$ be the intended smooth projective model over $\Q$ of the affine surface
\[y^2=g(t)(x^2+1),\qquad z^2=g(t)(x^2+2).\]
The source describes this as a one-parameter genus-one family and attributes to Skorobogatov $S(\Q)=\varnothing$ but $S(\A_{\Q})^{\operatorname{Br}}\ne\varnothing$. The Brauer--Manin set consists of adelic points orthogonal to the Brauer group under the sum of local invariants.

\textbf{Original question (Colliot-Thélène).} Is $S(k)\ne\varnothing$ for some number field $k/\Q$ of odd degree? The original remarks motivate this by existence of a zero-cycle of degree one and assert the converse implication for this example. A zero-cycle is an integral formal sum of closed points, with degree the corresponding sum of residue degrees; it need not be effective and does not itself give a rational point. The implication from an odd-degree point to degree-one index uses additional coprime-degree cycle information for this particular surface, not a general principle for all varieties.

The remarks say random searching even in a fixed cubic field is unlikely to find a point; propose a quotient of $S\times S\times S$ with an unspecified ``action by three'' (Ellenberg); suggest fixing $x$ to obtain $y^2=g(t)(x^2+1)$ and then searching for $z$ (Voloch); and observe that $(y/z)^2=(x^2+1)/(x^2+2)$ wherever defined (Poonen). McCallum asks whether degree three should suffice, and Colliot-Thélène says not necessarily. Poonen proposes instead finding a cubic point on $w^2=(x^2+1)(x^2+2)$, putting it in Weierstrass form and searching for $t$.

\textbf{Attached unreviewed solution attribution.} The aid credits Brendan Creutz with an affirmative solution in 2017: the surface has a point over $L=\Q(\theta)$, where $\theta^3+\theta+1=0$, so degree three does suffice. It says it translates and exactly verifies Creutz's published point in the AIM convention, but the full $(x,t,y,z)$ coordinates are not in this supplied summary. The original negative expectations are historical, not a denial of this later claimed result.

The aid supplies a descent bridge from
\[D_0:\quad w^2=(x^2+1)(x^2+2)\]
to $E:v^2=u^3-6u^2+u$ by
\[u=2w+2x^2+3,\qquad v=2xu.\]
On the associated untwisted cover $Y^2=x^2+1$, $Z^2=x^2+2$, $w=YZ$, the square root of $u$ is $Y+Z$. The aid also lists a low-height twist certificate
\[a=9-4\theta+6\theta^2,\quad U_0=2133-1224\theta+1494\theta^2,\quad Y=\theta^2,\quad Z=1-\theta^2,\quad N_{L/\Q}(a)=1,\]
and says the Weierstrass coordinate is $u=a$. The twist model, role of $U_0$ and map to the AIM coordinates are absent. These displayed $Y,Z$ have sum one, so they cannot be substituted into the untwisted formula $u=(Y+Z)^2$ to obtain $a$ without the missing twist normalization. The two coordinate conventions are therefore kept separate. The full solution and cubic point remain credited to Creutz; no independent verification or new solution is claimed.
\subsection{Short English statement}
The problem asks whether a specified Brauer--Manin counterexample surface acquires a point over an odd-degree field. The aid credits Creutz's 2017 cubic-field solution and gives a descent bridge and twist data. The full point and twist normalization are missing, so the listed coordinates are not silently combined into an unsupported explicit point.
\subsection{Sources}
\begin{itemize}
\item Local UnsolvedMath registry, permanent ID 20000453, source code catalogue entry; \texttt{aim-algebraic-number-theory-notes.json}, notes[104], tag \texttt{question}.
\item AIM workshop \emph{Rational and integral points on higher dimensional varieties}, Question 20: \url{https://aimath.org/WWN/qptsurface2/qptsurface2.pdf}. The construction is attributed to Skorobogatov, question to Colliot-Thélène; remarks name Ellenberg, Voloch, Poonen and McCallum. Supplied verification clarifies the smooth projective model and powers; no PDF was consulted.
\item Attached research attempt 1 credits Brendan Creutz's 2017 published solution, without title, identifier or theorem locator. Result \texttt{full\_solution}, status \texttt{solved}, rigor label \texttt{proved}, novelty confidence low. It is machine generated, not peer reviewed, and requires independent verification before citation. Original literature field empty.
\end{itemize}
\end{document}
